\section{Conclusion}
\label{sec:conclusion}

We proposed a decentralized incremental Fuzzy PI controller for
bittide and connected its complete rule law to native PI through
explicit, implementation-checked calibration. The exact central
nonlinear map explains how the controller reshapes correction
increments according to the relationship between occupancy error
and its change. For the stated synchronous zero-delay fluid reference,
invariant-domain and Lyapunov analysis establish convergence of buffer
deviations and pairwise frequency differences, with positive physical
frequencies and the stated safety and execution constraints. A certified
trajectory prefix and contracting tail prove a strict headroom benefit
for one prescribed matched-gain initialization: Fuzzy excursion is
below 17.315 frames, while PI attains approximately 23.403701 frames.

Continuous-phase Callisto tests separately evaluate delayed local
execution. Matched-gain startup experiments show smaller measured
buffer excursions at three prescribed initial mismatches. After
equal-budget tuning and frozen selection, 186 untouched tests give
mean paired RMS improvements of 16.78\% and 7.85\% for the selected
random-priority Fuzzy controller in two bounded-random regimes,
alongside lower mean command activity. A distinct burst-priority
representative also improves IAE in the larger held-out burst.
These benefits have explicit operating limits and a common-frequency
shift cost; they establish conditional value rather than performance
ordering across all delivery conditions.

The contribution is an explicit nonlinear bittide controller with
both a rigorous reference-model margin result and independent
empirical evidence of useful synchronization and resource tradeoffs.
Extending the certificate to asynchronous delayed execution and
testing broader networks and physical links are the next steps
towards deployment under application-specific timing requirements.


% Previous placeholder, preserved before the completed section was integrated.
% \section{Conclusion}
% \label{sec:conclusion}
% \noindent\textbf{Draft status: to be written.}
% Complete this section after the evaluation and discussion are integrated,
% keeping reference-model guarantees distinct from measured Callisto results.
